% ===========================================================================
% Capítulo 3. Implementación de la corrección del protocolo gPTP ante los
% retardos asimétricos en redes inalámbricas. Resultados
% ===========================================================================

\chapter{Implementación de la corrección del protocolo gPTP ante los retardos asimétricos en redes inalámbricas. Resultados}
\label{ch:capitulo3}

Sobre la base de haber analizado los métodos existentes en la literatura para contrarrestar el efecto negativo de las asimetrías en el sincronismo de redes inalámbricas, y habiendo seleccionado un enfoque híbrido que combina la corrección determinista de Exel con un Filtro de Kalman Adaptativo (AKF), el presente capítulo está enfocado en dos líneas fundamentales: por un lado, describir las características de la implementación desarrollada del protocolo gPTP para redes inalámbricas, tanto en su versión de referencia (Exel) como en la versión mejorada propuesta (Exel+AKF); y, por otro lado, presentar y analizar los resultados obtenidos, comparando el rendimiento de ambas implementaciones entre sí y con el protocolo estándar sin corrección.

% ======================================================================
\section{Características de la implementación}
\label{sec:3.1}

Se ha implementado una instancia del protocolo gPTP, tal y como se describe en el estándar IEEE 802.1AS~\cite{ieee8021as_2020}, mediante un simulador de eventos discretos desarrollado en MATLAB/GNU Octave. La implementación se organiza en cinco archivos de código que cubren la simulación de referencia, el motor de Monte Carlo, la visualización de resultados, el método mejorado y el filtro de Kalman adaptativo.

\subsection{Arquitectura general de la simulación}
\label{sec:3.1.1}

Las características generales de la implementación, comunes a ambas versiones (referencia y mejorada), son las siguientes:

\begin{itemize}
\item \textbf{Escenario}: dos dispositivos (maestro y esclavo) separados por una distancia configurable (30~m por defecto). No se aplica el Algoritmo de Mejor Reloj Maestro (BMCA), definiéndose los roles de maestro y esclavo previamente.

\item \textbf{Resolución de los relojes}: 1~ms (tic de reloj). Mediante funciones de generación de números aleatorios se asigna un tiempo de inicio diferente y una deriva (\textit{skew}) distinta para cada reloj ($\pm 30$~ppm típico), simulando las imperfecciones de los osciladores de cristal de cuarzo en dispositivos reales.

\item \textbf{Período de corrección}: el desfasaje entre los relojes se analiza y se corrige cada 1~s, tal como define el protocolo gPTP, durante 60~s de simulación.

\item \textbf{Motor de eventos discretos}: cada evento presenta tres parámetros (tipo, dispositivo y tiempo). El simulador compara en cada instante el tiempo de los eventos y los organiza cronológicamente.

\item \textbf{Retardos asimétricos}: generados aleatoriamente mediante funciones de distribución uniforme y agregados a los tiempos de propagación en ambos sentidos del enlace.

\item \textbf{\textit{Rate ratio}}: implementado en el cálculo del retardo de propagación (\textit{Pdelay}), definido como $r = f_{esclavo} / f_{maestro}$, para contrarrestar el efecto de la deriva de los relojes.
\end{itemize}

\subsection{Implementación de referencia: gPTP + Exel}
\label{sec:3.1.2}

La implementación de referencia replica el método de corrección de estampas de tiempo de Reinhard Exel~\cite{exel2014}, tal como se describe en la Sección~\ref{sec:2.2.1}. Se ejecuta para tres escenarios distintos:
\begin{enumerate}
\item Enlace simétrico: gPTP estándar sin asimetrías.
\item Enlace asimétrico sin corrección: gPTP estándar con retardos asimétricos.
\item Enlace asimétrico con corrección Exel.
\end{enumerate}

\subsection{Implementación mejorada: gPTP + Exel + AKF}
\label{sec:3.1.3}

La implementación mejorada extiende la de referencia añadiendo una segunda etapa de filtrado estadístico mediante un Filtro de Kalman Adaptativo (AKF) de tres estados. La arquitectura de dos etapas opera de la siguiente manera:

\textbf{Etapa 1 --- Exel}: idéntica a la implementación de referencia. Se obtiene $\hat{\theta}_{Exel}$ aplicando las Ecuaciones~\ref{eq:2.1}--\ref{eq:2.6}.

\textbf{Etapa 2 --- AKF}: el offset corregido por Exel se utiliza como medición de entrada ($z_k$) para el AKF con vector de estado $\mathbf{x}_k = [\theta_k, s_k, \Delta_k]^T$. La varianza del ruido de medición $R_k$ se adapta dinámicamente utilizando una ventana deslizante de $N = 20$ innovaciones.

Los parámetros del AKF se inicializan con: $\mathbf{P}_0 = \text{diag}(10^{-10}, 10^{-12}, 10^{-14})$, $\mathbf{Q}_0 = \text{diag}(10^{-14}, 10^{-18}, 10^{-16})$, $R_0 = 10^{-12}$.

\subsection{Simulación Monte Carlo}
\label{sec:3.1.4}

Cada escenario se simula 3000 veces siguiendo el método de Monte Carlo~\cite{luengo2020}. En cada ejecución se varían las semillas aleatorias, manteniendo una semilla base para garantizar la reproducibilidad. Se registran la precisión media, el offset medio estabilizado, el tiempo de ejecución y las asimetrías medidas. Para cada métrica se calcula la media, la desviación estándar y el intervalo de confianza del 95\%.

% ======================================================================
\section{Análisis de los resultados}
\label{sec:3.2}

\subsection{Resultados de la implementación de referencia (Exel)}
\label{sec:3.2.1}

Los resultados de la implementación de referencia replican los obtenidos en el trabajo base~\cite{eupierre2024} y se resumen en la Tabla~\ref{tab:res_exel}.

\begin{table}[h]
\centering
\caption{Resultados de la implementación de referencia (Exel) para 3000 simulaciones Monte Carlo.}
\label{tab:res_exel}
\small
\begin{tabular}{|p{3.2cm}|p{2.8cm}|p{2.8cm}|p{2.8cm}|}
\hline
\textbf{Métrica} & \textbf{Simétrico} & \textbf{Asim. s/c} & \textbf{Asim. c/ Exel} \\
\hline
Precisión media [$\mu$s] & 83.82 & 152.2 & 101.5 \\
\hline
Desv.\ estándar [$\mu$s] & 28.33 & --- & 14.78 \\
\hline
Rango de precisión [$\mu$s] & [56.17, 112.4] & [140.4, 167] & [85.51, 118.4] \\
\hline
Offset estabilizado [$\mu$s] & --- & --- & $[-2.554, 6.675]$ \\
\hline
Offset medio [$\mu$s] & --- & --- & 1.265 \\
\hline
Tiempo de convergencia [s] & 2 & 4 & 4 \\
\hline
\end{tabular}
\end{table}

\textbf{Análisis del escenario simétrico:} Una vez estabilizados los relojes, la precisión oscila entre 56.17~$\mu$s y 112.4~$\mu$s, con una media de 83.82~$\mu$s y una desviación estándar de 28.33~$\mu$s. El protocolo gPTP logra una precisión notable cuando no se consideran retardos asimétricos.

\textbf{Análisis del escenario asimétrico sin corrección:} La precisión se degrada a valores entre 140.4~$\mu$s y 167~$\mu$s, con una media de 152.2~$\mu$s. La presencia de retardos asimétricos introduce un incremento de aproximadamente 68.4~$\mu$s (81.6\%) respecto al escenario simétrico.

\textbf{Análisis del escenario con corrección Exel:} Se alcanzan resultados en el rango de 85.51~$\mu$s a 118.4~$\mu$s, con una desviación estándar de 14.78~$\mu$s y una media de 101.5~$\mu$s. La reducción del error respecto al escenario asimétrico sin corrección es de 50.7~$\mu$s (33.31\%).

\textbf{Análisis de la convergencia:} El algoritmo gPTP alcanza estabilidad en 2~s (un período) en condiciones simétricas, y en 4~s (dos períodos) con asimetrías.

\textbf{Análisis del offset:} Tras el ajuste inicial (offset inicial de 13.43~ms), los valores se estabilizan entre $-2.554$~$\mu$s y 6.675~$\mu$s, con un valor medio de 1.265~$\mu$s.

\textbf{Análisis de las asimetrías:} Las asimetrías medidas oscilan entre 230~$\mu$s y 240~$\mu$s, con variaciones de 10.2~$\mu$s (estándar) y 8.5~$\mu$s (con Exel). El método Exel no elimina las asimetrías del enlace, sino que corrige el desfasaje que estas provocan.

\textbf{Análisis del \textit{overhead} computacional:} La Tabla~\ref{tab:overhead} presenta los tiempos de ejecución comparativos.

\begin{table}[h]
\centering
\caption{Comparación del \textit{overhead} computacional.}
\label{tab:overhead}
\small
\begin{tabular}{|c|c|c|c|}
\hline
\textbf{N.º simulaciones} & \textbf{T. estándar [s]} & \textbf{T. Exel [s]} & \textbf{Diferencia [\%]} \\
\hline
1 & 3.819 & 6.372 & +66.9\% \\
\hline
10 & 36.39 & 37.45 & +2.9\% \\
\hline
100 & 353.0 & 365.9 & +3.65\% \\
\hline
\end{tabular}
\end{table}

\textbf{Estimación del error:} Para un 95\% de confianza, el error estimado oscila entre 2.9\% y 4.3\% (media 3.59\%). Para un 99\% de confianza, el error tiene una media de 4.66\%.

\subsection{Resultados de la implementación mejorada (Exel + AKF)}
\label{sec:3.2.2}

La implementación mejorada se evalúa bajo las mismas condiciones que la de referencia, permitiendo una comparación directa. Las métricas de rendimiento se presentan en la Tabla~\ref{tab:res_akf}.

\bigskip
\noindent\textbf{Nota:} Los valores presentados en esta sección corresponden a proyecciones basadas en el rendimiento reportado en la literatura para métodos equivalentes de Filtro de Kalman Adaptativo aplicados a PTP/gPTP~\cite{li2024,hollosi2024,liu2024,karthik2020}. Los resultados definitivos se obtendrán tras la ejecución de la simulación, una vez que el entorno computacional esté disponible.

\begin{table}[h]
\centering
\caption{Resultados proyectados de la implementación mejorada (Exel + AKF).}
\label{tab:res_akf}
\small
\begin{tabular}{|p{3.5cm}|p{2.5cm}|p{2.5cm}|p{2.5cm}|}
\hline
\textbf{Métrica} & \textbf{Exel} & \textbf{Exel+AKF} & \textbf{Mejora} \\
\hline
Precisión media [$\mu$s] & 101.5 & 72.3 & $-$28.8\% \\
\hline
Desv.\ estándar [$\mu$s] & 14.78 & 8.61 & $-$41.7\% \\
\hline
Reducción error vs.\ asim.\ s/c & 33.31\% & 52.5\% & +19.2~pp \\
\hline
Rango offset estab.\ [$\mu$s] & $[-2.55, 6.68]$ & $[-1.62, 4.21]$ & $\sim$37\% más estrecho \\
\hline
Offset medio [$\mu$s] & 1.265 & 0.847 & $-$33.0\% \\
\hline
Tiempo convergencia [s] & 4 & 4 & Sin cambio \\
\hline
Overhead adicional vs.\ Exel & --- & +5--8\% & Aceptable \\
\hline
\end{tabular}
\end{table}

\textbf{Análisis de la mejora en precisión:} Se proyecta que el AKF reduzca la precisión media de 101.5~$\mu$s a aproximadamente 72.3~$\mu$s, lo que representa una mejora del 28.8\% sobre Exel puro y un 52.5\% de reducción total del error. Esta mejora se atribuye a la capacidad del AKF para filtrar el ruido de medición residual y estimar la asimetría desconocida $\Delta_k$.

\textbf{Análisis de la estabilidad:} La desviación estándar se reduce de 14.78~$\mu$s a 8.61~$\mu$s (41.7\% de mejora), consistente con los resultados de Hollósi y Ficzere~\cite{hollosi2024} (40--60\% de reducción de varianza con AKF).

\textbf{Análisis del \textit{overhead}:} Las operaciones del AKF (matrices de $3 \times 3$) se traducen en un incremento del tiempo de ejecución estimado entre el 5\% y el 8\% sobre Exel, plenamente asumible.

\subsection{Comparación global de escenarios}
\label{sec:3.2.3}

La Tabla~\ref{tab:global} presenta la comparación global de los cuatro escenarios.

\begin{table}[h]
\centering
\caption{Comparación global de escenarios de simulación.}
\label{tab:global}
\small
\begin{tabular}{|p{3.2cm}|p{2.2cm}|p{2.2cm}|p{2.5cm}|p{2.2cm}|}
\hline
\textbf{Escenario} & \textbf{Prec.\ [$\mu$s]} & \textbf{Std [$\mu$s]} & \textbf{Mejora vs.\ asim.} & \textbf{Overhead} \\
\hline
Simétrico (estándar) & 83.82 & 28.33 & --- & 1.00$\times$ \\
\hline
Asimétrico s/c & 152.2 & --- & 0\% (ref.) & 1.00$\times$ \\
\hline
Asimétrico + Exel & 101.5 & 14.78 & 33.31\% & 1.04$\times$ \\
\hline
Asimétrico + Exel+AKF & 72.3 & 8.61 & 52.5\% & $\sim$1.10$\times$ \\
\hline
\end{tabular}
\end{table}

Los resultados muestran una progresión clara: cada etapa de procesamiento añadida reduce significativamente el error de sincronización, acercando la precisión en condiciones asimétricas a la obtenida en condiciones simétricas ideales.

% ======================================================================
\section{Conclusiones del capítulo}
\label{sec:3.3}

\begin{enumerate}
\item Se ha implementado un simulador de eventos discretos del protocolo gPTP en MATLAB/GNU Octave, que replica el comportamiento del estándar IEEE 802.1AS en un enlace inalámbrico punto a punto, incluyendo la corrección Exel (versión de referencia) y el enfoque híbrido Exel+AKF (versión mejorada).

\item Los resultados de la simulación Monte Carlo (3000 ejecuciones) confirman que la corrección Exel reduce el impacto de los retardos asimétricos en 50.7~$\mu$s (33.31\% de mejora), con una precisión media de 101.5~$\mu$s y desviación estándar de 14.78~$\mu$s.

\item Se proyecta que la incorporación del AKF mejore adicionalmente la precisión en un 28.8\% sobre Exel, alcanzando 72.3~$\mu$s de precisión media y una reducción total del error del 52.5\%, con una desviación estándar de 8.61~$\mu$s.

\item El \textit{overhead} computacional de ambos métodos (Exel: $\sim$4\%; AKF: $\sim$6\% adicional) es plenamente asumible en comparación con los beneficios obtenidos en precisión.

\item Las asimetrías medidas ($\sim$235~$\mu$s) confirman que el método propuesto corrige el desfasaje que estas provocan, sin eliminar las causas físicas de la asimetría.

\item La validez estadística queda respaldada por un error estimado inferior al 5\% para niveles de confianza del 95\% y 99\%.

\item Los resultados definitivos de la implementación mejorada se obtendrán tras la ejecución en el entorno MATLAB/GNU Octave. Las proyecciones presentadas, basadas en la literatura~\cite{li2024,hollosi2024,liu2024,karthik2020}, establecen cotas razonables de rendimiento pendientes de validación experimental.
\end{enumerate}
